% Part: propositional-logic
% Chapter: syntax-and-semantics
% Section: valuations-sat

\documentclass[../../../include/open-logic-section]{subfiles}

\begin{document}

\olfileid{pl}{syn}{val}

\olsection{\usetoken{P}{valuation} మరియు సంతృప్తి}

\begin{defn}[!!^{valuation}s]
“సత్యం,” “అసత్యం” అనే రెండు సత్యమూల్యాల సమితి
$\{\True, \False\}$ అనుకుందాం. $\Lang{L_0}$కు ఒక
\emph{!!{valuation}} అంటే భాషలోని !!{propositional variable}sకు
$\True$ లేదా $\False$ను కేటాయించే ప్రమేయం~$\pAssign{v}$; అంటే
$\pAssign{v} \colon \PVar \to \{\True, \False \}$.
\end{defn}

\begin{defn}
  ఒక !!a{valuation}~$\pAssign{v}$ ఇచ్చినప్పుడు, మూల్యాంకన ప్రమేయం
  $\pValue{v} \colon \Frm[L_0] \to \{\True, \False \}$ను కింది
  విధంగా ఆగమనాత్మకంగా నిర్వచించండి:
  \begin{align*}
    \iftag{prvFalse}{\pValue{v}(\lfalse) & = \False; \\}{}
    \iftag{prvTrue}{\pValue{v}(\ltrue) & = \True; \\}{}
    \pValue{v}(\Obj p_n) & = \pAssign{v}(\Obj p_n); \\
    \iftag{prvNot}{\pValue{v}(\lnot !A) & = \begin{cases}
        \True & \text{if } \pValue{v}(!A) = \False;\\
        \False & \text{otherwise.}
      \end{cases} \\ }{}
    \iftag{prvAnd}{\pValue{v}(!A \land !B) & = \begin{cases}
        \True &
        \text{if $\pValue{v}(!A) = \True$ and $\pValue{v}(!B) = \True$;}\\
        \False &
        \text{if $\pValue{v}(!A) = \False$ or $\pValue{v}(!B) = \False$}.
      \end{cases}\\}{}
    \iftag{prvOr}{\pValue{v}(!A \lor !B) & = \begin{cases}
        \True &
        \text{if $\pValue{v}(!A) = \True$ or $\pValue{v}(!B) = \True$;}\\
        \False &
        \text{if $\pValue{v}(!A) = \False$ and $\pValue{v}(!B) = \False$}.
      \end{cases}\\}{}
    \iftag{prvIf}{\pValue{v}(!A \lif !B) & = \begin{cases}
        \True &
        \text{if $\pValue{v}(!A) = \False$ or $\pValue{v}(!B) = \True$;}\\
        \False &
        \text{if $\pValue{v}(!A) = \True$ and $\pValue{v}(!B) = \False$}.
      \end{cases}\\}{}
    \iftag{prvIff}{\pValue{v}(!A \liff !B) & = \begin{cases}
        \True &
        \text{if $\pValue{v}(!A) = \pValue{v}(!B)$;}\\
        \False &
        \text{if $\pValue{v}(!A) \neq \pValue{v}(!B)$}.
      \end{cases}}{}
\end{align*}
\end{defn}

\begin{explain}
ఈ ఉపవాక్యాలు కింది సత్య పట్టికలకు సరిపోతాయి:
\begin{center}
\iftag{prvNot}{
  \begin{tabular}{|c||c|} \hline
    $!A$ & $ \lnot !A$ \\
    \hline \hline
    $\True$ & $\False$ \\
    $\False$ & $\True$ \\
    \hline
  \end{tabular}
}{}
\iftag{prvAnd}{
  \begin{tabular}{|cc||c|} \hline
    $!A$ & $!B$ & $!A \land !B$ \\
    \hline \hline
    $\True$ & $\True$ & $\True$ \\
    $\True$ & $\False$ & $\False$ \\
    $\False$ & $\True$ & $\False$ \\
    $\False$ & $\False$ & $\False$ \\
    \hline
  \end{tabular}
}{}
\iftag{prvOr}{
  \begin{tabular}{|cc||c|} \hline
    $!A$ & $!B$ & $!A \lor !B$ \\
    \hline \hline
    $\True$ & $\True$ & $\True$ \\
    $\True$ & $\False$ & $\True$ \\
    $\False$ & $\True$ & $\True$ \\
    $\False$ & $\False$ & $\False$ \\
    \hline
  \end{tabular}
}{}%
\iftag{notprvNot,notprvAnd,notprvOr,notprvIf}{}{\\[1em]}
\iftag{prvIf}{
  \begin{tabular}{|cc||c|} \hline
    $!A$ & $!B$ & $!A \lif !B$ \\
    \hline \hline
    $\True$ & $\True$ & $\True$ \\
    $\True$ & $\False$ & $\False$ \\
    $\False$ & $\True$ & $\True$ \\
    $\False$ & $\False$ & $\True$ \\
    \hline
  \end{tabular}
}{}
\iftag{prvIff}{
  \begin{tabular}{|cc||c|} \hline
    $!A$ & $!B$ & $!A \liff !B$ \\
    \hline \hline
    $\True$ & $\True$ & $\True$ \\
    $\True$ & $\False$ & $\False$ \\
    $\False$ & $\True$ & $\False$ \\
    $\False$ & $\False$ & $\True$ \\
    \hline
  \end{tabular}
}{}
\end{center}
\end{explain}

\begin{prob}
$\Lang{L_0}$కు $\diamondsuit$ అనే త్రిస్థాన సంయోజకాన్ని చేర్చడాన్ని
పరిశీలించండి; దాని మూల్యాంకనం ఇలా ఇవ్వబడింది:
\begin{gather*}
  \pValue{v}(\diamondsuit( !A , !B , !C ) ) = \begin{cases}
    \pValue{v}( !B ) &
    \text{if $\pValue{v}(!A) = \True$;}\\
    \pValue{v}( !C ) &
    \text{if $\pValue{v}(!A) = \False $}.
  \end{cases}
\end{gather*}
ఈ సంయోజకానికి సత్య పట్టికను రాయండి.
\end{prob}

\begin{thm}[స్థానిక నిర్ణయితత్వం]
\ollabel{thm:LocalDetermination} $\pAssign{v_1}$, $\pAssign{v_2}$ అనే
!!{valuation}s, $!A$లో సంభవించే !!{propositional
variable}sపై
ఏకీభవిస్తాయని అనుకుందాం; అంటే, ఆ నిర్దిష్ట
!!{formula}~$!A$లో $\Obj p_n$ సంభవించినప్పుడల్లా
$\pAssign{v_1}(\Obj p_n) = \pAssign{v_2}(\Obj p_n)$. అప్పుడు
$\pValue{v_1}$, $\pValue{v_2}$లు $!A$పై కూడా ఏకీభవిస్తాయి; అంటే
$\pValue{v_1}(!A) = \pValue{v_2}(!A)$.
\end{thm}
\sourcecorrection{OLTEPLSYN-004}{ఆంగ్ల మూలం మొదట స్థిరపరిచిన $!A$లో
సంభవించే చరరాశుల గురించి చెప్పి, వివరణలో “ఏదో సూత్రం $!A$లో” అని
రాసింది. అస్తిత్వ పరిమాణీకరణగా చదివితే అది స్థానిక పరిమితిని
కోల్పోతుంది. సిద్ధాంతపు ముందు వాక్యం, ముగింపు నిర్ణయించినట్లుగా
“ఆ నిర్దిష్ట సూత్రం $!A$లో” అని స్పష్టం చేశాం.}

\begin{proof}
$!A$పై ఆగమనం ద్వారా.
\end{proof}

\begin{defn}[సంతృప్తి]
\ollabel{defn:satisfaction} \emph{ఒక !!a{formula}~$!A$ను
  !!a{valuation}~$\pAssign{v}$ సంతృప్తిపరచడం} అనే భావన
  $\pSat{v}{!A}$ను కింది విధంగా ఆగమనాత్మకంగా నిర్వచించవచ్చు.
  (“$\pSat{v}{!A}$ కాదు” అని సూచించడానికి $\pSat/{v}{!A}$ అని
  రాస్తాం.)
\begin{enumerate}
\tagitem{prvFalse}{%
  \indcase{!A}{\lfalse}{$\pSat/{v}{\indfrm}$.}}{}

\tagitem{prvTrue}{%
  \indcase{!A}{\ltrue}{$\pSat{v}{\indfrm}$.}}{}

\item \indcase{!A}{\Obj p_i}{$\pSat{v}{\indfrm}$ కావడం,
  $\pAssign{v}(\Obj p_i) = \True$ అయినప్పుడూ, అప్పుడు మాత్రమే.}

\tagitem{prvNot}{%
  \indcase{!A}{\lnot !B}{$\pSat{v}{\indfrm}$ కావడం,
    $\pSat/{v}{!B}$ అయినప్పుడూ, అప్పుడు మాత్రమే.}}{}

\tagitem{prvAnd}{%
  \indcase{!A}{(!B \land !C)}{$\pSat{v}{\indfrm}$ కావడం,
    $\pSat{v}{!B}$ మరియు $\pSat{v}{!C}$ అయినప్పుడూ, అప్పుడు మాత్రమే.}}{}

\tagitem{prvOr}{%
  \indcase{!A}{(!B \lor !C)}{$\pSat{v}{\indfrm}$ కావడం,
    $\pSat{v}{!B}$ లేదా $\pSat{v}{!C}$ (లేదా రెండూ) అయినప్పుడూ,
    అప్పుడు మాత్రమే.}}{}

\tagitem{prvIf}{%
  \indcase{!A}{(!B \lif !C)}{$\pSat{v}{\indfrm}$ కావడం,
    $\pSat/{v}{!B}$ లేదా $\pSat{v}{!C}$ (లేదా రెండూ) అయినప్పుడూ,
    అప్పుడు మాత్రమే.}}{}

\tagitem{prvIff}{%
  \indcase{!A}{(!B \liff !C)}{$\pSat{v}{\indfrm}$ కావడం,
    $\pSat{v}{!B}$, $\pSat{v}{!C}$ రెండూ, లేదా
    $\pSat{v}{!B}$, $\pSat{v}{!C}$ రెండూ కాకపోవడం జరిగినప్పుడూ,
    అప్పుడు మాత్రమే.}}{}
\end{enumerate}
$\Gamma$ ఒక !!{formula}s సమితి అయితే, $\pSat{v}{\Gamma}$ కావడం,
ప్రతి~$!A \in \Gamma$కు $\pSat{v}{!A}$ అయినప్పుడూ, అప్పుడు మాత్రమే.
\end{defn}

\begin{prop}\ollabel{prop:sat-value}
  $\pSat{v}{!A}$ కావడం, $\pValue{v}(!A) = \True$ అయినప్పుడూ,
  అప్పుడు మాత్రమే.
\end{prop}

\begin{proof}
  $!A$పై ఆగమనం ద్వారా.
\end{proof}

\begin{prob}
  \olref[pl][syn][val]{prop:sat-value}ను నిరూపించండి.
\end{prob}

\end{document}
